%!TEX program = lualatex
% ============================================================
%  国家自然科学基金（NSFC）报告正文 LaTeX 模板  (2026 版)
%  Author: Lily Xue  (https://github.com/<your-handle>)
%  License: MIT
%
%  编译命令（推荐 LuaLaTeX）：
%      lualatex main
%      bibtex   main
%      lualatex main
%      lualatex main
%
%  也可用 XeLaTeX：把第 1 行改为 %!TEX program = xelatex 即可。
% ============================================================

\documentclass[12pt]{article}
\usepackage[UTF8,fontset=fandol]{ctex}
\setmainfont[Ligatures=TeX]{Latin Modern Roman}

% 楷体字体（用于"提纲提示语"）
\setCJKfamilyfont{zhkai}[AutoFakeBold=5]{FandolKai-Regular}
\newcommand{\kai}[1]{{\kaishu #1}}
\newcommand{\kaibold}[1]{{\bfseries\kaishu #1}}

\usepackage{nsfc}            % 本模板核心样式
\usepackage[misc]{ifsym}
\usepackage{pifont}
\usepackage{overpic}
\usepackage{subfigure}
\usepackage{enumitem}
\usepackage{color}
\usepackage{multirow}
\usepackage{booktabs}
\usepackage{threeparttable}
\usepackage{textcomp}
\usepackage{float}

% ---------- 常用宏 ----------
\newcommand{\cmark}{\ding{51}}
\newcommand{\xmark}{\text{\ding{55}}}
\newcommand{\todo}[1]{{\textcolor{red}{\bf [#1]}}}
\newcommand{\mypara}[1]{\paragraph{#1.}}
\newcommand{\myEmph}[1]{\textbf{\textcolor[rgb]{0,0,0.25}{#1}}}
\newcommand{\addImg}[2]{\includegraphics[width=#1\textwidth]{#2}}

% ---------- 申请人标注 ----------
\newcommand{\me}[1]{{\CJKfamily{hei}#1}}                % 申请人用黑体
\newcommand{\first}{\textsuperscript{\textdagger}}      % 一作符号
\newcommand{\cor}{\textsuperscript{*}}                  % 通讯作者符号

% ---------- 引用辅助 ----------
\newcommand{\figref}[1]{图\ref{#1}}
\newcommand{\tabref}[1]{表\ref{#1}}
\newcommand{\equref}[1]{式\ref{#1}}
\newcommand{\secref}[1]{第\ref{#1}节}
\newcommand{\citess}[1]{\textsuperscript{\cite{#1}}}    % 上标引用
\newcommand{\myCite}[1]{\myEmph{#1}$^\text{我们}_\text{工作}$}  % 标注申请人自己的工作

% ---------- 压缩参考文献间距 ----------
\let\origthebibliography\thebibliography
\let\origendthebibliography\endthebibliography
\renewenvironment{thebibliography}[1]{%
  \origthebibliography{#1}%
  \setlength{\itemsep}{0pt}\setlength{\parsep}{0pt}%
  \setlength{\topsep}{0pt}\setlength{\parskip}{0pt}%
}{\origendthebibliography}

\WarningFilter{latexfont}{Font}
\WarningFilter{font}{Font}
\WarningFilter{fontspec}{Font}

\graphicspath{{figures/}}

\begin{document}

\title{\emph{\zihao{2} \kai{报告正文（2026版）}}}
\maketitle
\thispagestyle{empty}

\kai{参照以下提纲撰写，要求内容翔实、清晰，层次分明，标题突出。\textbf{\textcolor{red}{申请书正文原则上不超过30页，鼓励简洁表达。}\textcolor[rgb]{0.0625,0.4453,0.7383}{请勿删除或改动下述提纲标题及括号中的文字。}}}

\vspace{0.5cm}

% ============================================================
\ContentDes{\kaibold{（一）立项依据：}}
\NsfcNote{\kai{（为什么要开展此项研究，研究的科学技术价值如何）}}

\section{研究背景及科学意义}
% 写作建议：
%  1. 先用 1-2 段勾勒国家战略与现实痛点（政策文件 + 行业数据）
%  2. 再用 1-2 段梳理国际前沿（高水平期刊代表性工作）
%  3. 最后用 1 段引出本项目的切入点与研究范式
\todo{在这里写研究背景}

\subsection{关键科学问题的界定}
% 建议用 3 个并列子段落，每段聚焦一个科学问题：
%   (1) ...
%   (2) ...
%   (3) ...
\todo{列出 2-3 个关键科学问题}

\section{国内外研究现状与发展动态}
\todo{文献综述。建议按"科学问题"组织，而非按"作者"罗列。}

\section{申请人前期工作基础}
% 突出本团队过去 3-5 年与本课题最相关的产出
\todo{列出前期工作：已发表论文、在投论文、数据/方法积累}

% ============================================================
\ContentDes{\kaibold{（二）项目研究内容：}}
\NsfcNote{\kai{（围绕关键科学问题，开展哪些研究）}}

\section{研究目标}
\todo{1-2 段，明确"做什么 + 解决什么科学问题 + 形成什么理论或方法"}

\section{研究内容}
\subsection{子课题一：\todo{标题}}
\subsection{子课题二：\todo{标题}}
\subsection{子课题三：\todo{标题}}

\section{拟解决的关键科学问题}
\todo{呼应（一）中的关键科学问题，逐条说明本项目如何攻克}

% ============================================================
\ContentDes{\kaibold{（三）研究方案及可行性分析：}}
\NsfcNote{\kai{（怎么开展此项研究，开展研究的把握如何）}}

\section{研究方案}
\subsection{总体研究框架}
% 强烈建议放一张"研究框架图"
% \begin{figure}[H]
%   \centering
%   \addImg{0.85}{framework.png}
%   \caption{\kai{本项目研究框架}}
%   \label{fig:framework}
% \end{figure}
\todo{放研究框架图 + 1 段文字描述各子课题之间的逻辑递进}

\subsection{子课题一研究方案}
\subsection{子课题二研究方案}
\subsection{子课题三研究方案}

\section{可行性分析}
\subsection{理论与方法可行性}
\subsection{数据与平台可行性}
\subsection{团队与时间可行性}

% ============================================================
\ContentDes{\kaibold{（四）本项目的特色与创新之处：}}

\begin{itemize}
  \item \myEmph{特色一：} \todo{...}
  \item \myEmph{特色二：} \todo{...}
  \item \myEmph{创新一：} \todo{...}
  \item \myEmph{创新二：} \todo{...}
\end{itemize}

% ============================================================
\ContentDes{\kaibold{（五）年度研究计划及预期研究结果：}}

\subsection*{年度研究计划}
\begin{itemize}
  \item \textbf{2027.01--2027.12：} \todo{...}
  \item \textbf{2028.01--2028.12：} \todo{...}
  \item \textbf{2029.01--2029.12：} \todo{...}
  \item \textbf{2030.01--2030.12：} \todo{...}
\end{itemize}

\subsection*{预期研究结果}
\todo{列出预期论文（标注目标期刊）、专著、数据、政策报告、人才培养等}

% ============================================================
%  参考文献
% ============================================================
\bibliographystyle{nsfc}
\bibliography{references}

% ============================================================
%  研究基础与工作条件
% ============================================================
\newpage
\ContentDes{\kaibold{（一）研究基础}}
\NsfcNote{\kai{（与本项目相关的研究工作积累和已取得的研究工作成绩）}}
\todo{...}

\ContentDes{\kaibold{（二）工作条件}}
\NsfcNote{\kai{（包括已具备的实验条件，尚缺少的实验条件和拟解决的途径，包括利用国家实验室、国家重点实验室和部门重点实验室等研究基地的计划与落实情况）}}
\todo{...}

\ContentDes{\kaibold{（三）正在承担的与本项目相关的科研项目情况}}
\todo{...}

\ContentDes{\kaibold{（四）完成国家自然科学基金项目情况}}
\todo{...}

\end{document}
